On-device assistants must act on a user's behalf (setting an alarm, toggling a
setting, placing a call) without sending the request to the cloud. This thesis
asks how much language model such \emph{action} actually requires. We present
A.L.F.R.E.D., a local-first middleware that splits every request between a tiny
on-device \emph{reflexer} (a 2B model that fills a distilled command pattern in a
single forward pass) and a heavier \emph{thinker} (an agent loop) reserved for the
long tail. Across purpose-built benchmarks scored by an end-state oracle, four
results establish the central claim. First, on binding, the 2B reflexer with a
distilled pattern matches a $17\times$ larger free-form thinker (both 100\%),
while the same 2B without the pattern reaches only 30\%; distillation, not scale,
is what produces reliable execution. Second, the abstraction's benefit shifts with
the interface: on a clean settings API it is an efficiency win ($+65$ points at two
commands instead of nine), and on a realistic Android API it becomes an accuracy
win. Third, and most importantly, even a 35B reasoning model fails the device's
undiscoverable encoding conventions. Model scale buys \emph{discovery} (which key)
but not \emph{convention} (what the value means), and only a distilled pattern
supplies the convention. Fourth, that pattern can be minted automatically by a
teacher that never sees the tasks, lifting the 2B reflexer from 25\% to 76--80\%,
and the choice of teacher propagates directly to the student (a cloud teacher's
patterns score 80\% against a local teacher's 60\%). The remaining open limit is
routing, not execution. We frame these findings as the \emph{executioner
paradigm}: distill device conventions once with a strong teacher, then apply them
cheaply and repeatedly with a small on-device reflexer.

\vspace{0.6em}
\noindent\textbf{Keywords:} on-device language models, knowledge distillation,
agent routing, tool use, intent execution, privacy.
